\section{Edge-Device Latency \& Energy Methodology}
\label{app:latency-methodology}

This appendix documents the end-to-end protocol used to produce the
single-user, single-stream latency and energy numbers reported for the
\textsc{Spark} edge node (NVIDIA GB10, 119.7\,GiB unified memory, ARM CPU).
The goal is a faithful estimate of the latency a single end-user would
experience on a personal device, separating the cost of memory backend
mechanics from the cost of LLM inference.

\subsection{Hardware and serving stack}
\label{app:lat:hardware}

All measurements run on a single \textsc{Spark} GB10 node. The reader LLM is
served by \texttt{sglang} (\texttt{lmsysorg/sglang:spark}) inside Docker, with
\texttt{--max-running-requests 1} on the server and
\texttt{--answer-concurrency 1} on the client; this disables continuous
batching so each request is processed in isolation. Memory subsystems run as
peers on the same GPU: \textsc{Memobase} ships its full Postgres+server stack
via a per-trial \texttt{docker compose} bundle, and \textsc{MemSearch} uses
embedded \texttt{milvus-lite} plus a separately-loaded
\texttt{ollama-memos} container that serves \texttt{nomic-embed-text}.
Co-residency on a 119\,GiB unified-memory device forces an explicit
\texttt{--mem-fraction-static} budget for sglang
(\texttt{0.85} for vanilla/oracle/memobase, \texttt{0.5} when ollama coexists).
A 100\,ms-period \texttt{hw\_probe} sidecar samples GPU power, GPU/CPU
utilization, GPU/CPU temperature, and SoC power rails throughout every cell.

\paragraph{Failure modes encountered.} ollama was observed to silently
load the embedder onto CPU when the GPU was already partially occupied at
container start (a single embed call grew from $\sim$30\,ms to
$>$2\,min). The protocol therefore mandates that ollama be (re)started
\emph{before} sglang, with a verification step (a warm \texttt{/api/embed} that
must return in $<$1\,s and a docker-log assertion of
\texttt{offloaded 13/13 layers to GPU}) before any timed work begins.

\subsection{Query selection and namespace}
\label{app:lat:queries}

To make the measurement budget tractable while keeping the question pool
representative, we select $90$ queries via
\texttt{scripts/select\_latency\_queries.py}: nine sub-dimensions
($d_1$~conflict, $d_2$~anaphora, $d_3$~confabulation,
$d_4$~permission, $d_5$~cloze, $d_6$~metadata, $d_7$~qa,
$d_8$~temporal, $d_{10}$~counterfactual), $10$ latest
queries per sub-dimension by \texttt{metadata.query\_timestamp},
falling back to $\max(\text{evidence\_session.end\_time})$ where
\texttt{query\_timestamp} is absent (e.g.\ \DFour permission). The result
is written as a JSON list and consumed by \texttt{eval/cli.py} via the
\texttt{--instance-id-file} flag, restricting the QA pool to those
exact $90$ instance ids. Memory subsystems use a \emph{namespace per
ego} (\texttt{namespace\_mode=ego\_agent\_id}); ingest is monotonic
day-by-day so no future evidence leaks into earlier days.

\subsection{Phase 1 — direct measurement of reader-bound backends}
\label{app:lat:phase1}

Vanilla, Oracle, and In-Memory RAG (\texttt{inmem}) are evaluated as
$\{\text{reader}\}\times\{\text{backend}\}$ cells via
\texttt{scripts/run\_latency\_answer\_sweep.sh}. For each reader the
sweep brings up sglang with \texttt{--max-running-requests 1}, performs
a chat-completion sanity probe (a single ``Reply OK.'' request), and
only then begins timed work; an empty or malformed sanity output aborts
the cell to fail-fast on weight corruption. Each cell records, per
question: prompt tokens $N_p$, completion tokens $N_c$, time-to-first-token
$\textsc{ttft}$, total answer time $T_{\text{ans}}$, and a per-cell
\texttt{hw\_agg} JSON aggregating the \texttt{hw\_probe} CSV. The
\texttt{eval} configuration is
\texttt{eval/config/pipeline\_vanilla512.yaml} so the Vanilla budget is
held at $512$ session tokens, matching the headline ablation (commit
\texttt{7cb42d9}).

\subsection{Phase 2 — split ingest \& answer for cache-replay backends}
\label{app:lat:phase2}

\textsc{MemSearch} and \textsc{Memobase} expose no live answer-time
adapter in our pipeline; both produce a per-question memory cache during
\emph{ingest} that the answer phase then replays. We measure them in
two steps:

\begin{enumerate}
\item \textbf{Ingest.} \texttt{scripts/run\_latency\_memsearch\_ingest.sh}
and \texttt{run\_latency\_memobase\_ingest.sh} bring up sglang plus the
backend's storage stack, run \texttt{scripts/reproduce/build\_memory\_cache.py}
on $8$ ego agents $\times\,15$ days at concurrency $1$, and capture
\texttt{hw\_probe}. The cache JSONL written by
\texttt{build\_memory\_cache.py} stores, per question, the live retrieval
latency observed during ingest:
\begin{equation*}
\texttt{retrieve\_latency\_ms} \;\equiv\; T_{\text{search}}\!\big(\text{backend, query}\big).
\end{equation*}
Because retrieval is a deterministic function of (query, index) and is
disjoint from reader-LLM work, \emph{this column is the
reader-independent live search overhead}; the answer phase does not need
to repeat the search to expose it.

\item \textbf{Answer-replay.} \texttt{scripts/run\_latency\_answer.sh}
maps \texttt{BACKEND}$\in\{$\texttt{memsearch}, \texttt{memobase}$\}$
onto the eval-CLI's \texttt{--system memory\_cache} backend, replaying
each cached row through sglang to obtain $(N_p, N_c, T_{\text{ans}},
\textsc{ttft})$. Replay reports $T_{\text{search}}=0$ by construction,
which we replace with the per-row \texttt{retrieve\_latency\_ms} from
the ingest cache when fitting (\S\ref{app:lat:fit}).
\end{enumerate}

To avoid cache misses, each cell uses a backend-specific
\texttt{cache\_ids\_<backend>\_<reader>\_s2.json} filter listing the
$\sim\!430$ instance ids actually present in that backend's cache.
Because the cache is built on $8$ egos but the curated $90$-query
selection spans all $50$ egos, the in-cache subset is what actually
gets timed.

\subsection{Composition: reader-bound LLM time + reader-free search}
\label{app:lat:composition}

The single decomposition that makes a $5\!\times\!5$ table tractable
without running every cell is
\begin{equation}
T_{\text{total}}(\text{reader}, \text{backend}, q) \;=\;
\underbrace{T_{\text{LLM}}(\text{reader},\, N_p^q, N_c^q)}_{\text{reader-bound, backend-orthogonal}}
\;+\;
\underbrace{T_{\text{search}}(\text{backend})}_{\text{reader-orthogonal}}.
\label{eq:lat-composition}
\end{equation}
The first term depends on the reader and the per-question prompt /
completion lengths; the second is set by the memory backend's retrieval
pipeline (BM25, embedding lookup, etc.) and the storage stack. This
factorization is exact for backends whose retrieval does not call the
reader (\textsc{Vanilla}, \textsc{Oracle}, \textsc{RAG},
\textsc{MemSearch}); it is approximate for \textsc{Memobase}, whose
\emph{ingest} uses the reader as extractor (the cache is therefore
extractor-specific), but whose \emph{retrieve} operation at answer time
is a Postgres call and reader-orthogonal.

\subsection{Per-reader LLM-gen fit}
\label{app:lat:fit}

Within Eq.~\ref{eq:lat-composition} we model
$T_{\text{LLM}}$ as linear in the per-question prompt and completion
length:
\begin{equation}
T_{\text{LLM}}(\text{reader},\, N_p, N_c) \;=\;
\alpha_r + \beta_r\, N_p + \gamma_r\, N_c.
\label{eq:lat-fit}
\end{equation}
Per-reader $(\alpha_r,\beta_r,\gamma_r)$ are obtained by ordinary
least-squares from \texttt{scripts/fit\_latency\_model.py} using
\emph{only} Phase-1 backends (vanilla / oracle / inmem) — these provide
$N_p \in [\sim\!200,\sim\!11{,}000]$ and $N_c \in [1, 500]$, spanning the
ranges Phase-2 cells will land in. We do not include
\textsc{MemSearch} / \textsc{Memobase} samples in the regression to keep
the LLM-time fit free of any extractor- or replay-induced confounds.

The $T_{\text{search}}$ term in Eq.~\ref{eq:lat-composition} is computed
separately, per backend, as the median over per-question search times:
$0$ for \textsc{Vanilla} / \textsc{Oracle}, the answer-phase
\texttt{search\_time\_ms} for \textsc{RAG}, and the cache's
\texttt{retrieve\_latency\_ms} for \textsc{MemSearch} / \textsc{Memobase}.

\begin{table}[!htbp]
\centering
\small
\caption{Per-reader linear LLM-generation latency fit on Spark GB10 (\texttt{concurrency=1}). The model is $T_{\text{LLM}} = \alpha + \beta N_p + \gamma N_c$. Fit data is the union of vanilla / oracle / inmem cells (\S\ref{app:lat:fit}). $|r|_{\text{p95}}$ is the absolute residual at the 95th percentile.}
\label{tab:appendix-latency-fit}
\begin{tabular}{lrrrrrr}
\toprule
Reader & $\alpha$ (ms) & $\beta$ (ms\,/\,1k prompt tok) & $\gamma$ (ms\,/\,compl tok) & $R^2$ & $n_{\text{obs}}$ & $|r|_{\text{p95}}$ (ms) \\
\midrule
Qwen3-0.6B & -74.5 & 48.05 & 10.54 & 0.713 & 270 & 163 \\
Llama-3.2-3B & -60.2 & 102.88 & 38.71 & 0.983 & 270 & 491 \\
Mistral-7B-Inst.\ v0.3 & -- & -- & -- & -- & -- & -- \\
Qwen3-8B & -167.1 & 216.65 & 78.67 & 0.925 & 270 & 1118 \\
Qwen3-32B-AWQ & -650.9 & 754.38 & 99.49 & 0.780 & 270 & 4106 \\
\bottomrule
\end{tabular}
\end{table}


\paragraph{TTFT.} The same script also fits time-to-first-token by
ordinary least squares,
\begin{equation}
T_{\text{TTFT,LLM}}(\text{reader},\, N_p) \;=\; \alpha'_r + \beta'_r\, N_p,
\label{eq:lat-fit-ttft}
\end{equation}
omitting the completion-token term (TTFT is prefill-only). End-to-end
TTFT for the main-text efficiency table is then
$T_{\text{TTFT}} = T_{\text{search}}(\text{backend}) + T_{\text{TTFT,LLM}}$,
with $T_{\text{search}}$ from the same per-backend column used in
\S\ref{app:lat:fit}.

\paragraph{Diagnostics.} The fit script reports per-reader $R^2$,
$|residual|$ at p50 and p95, observed $(N_p, N_c)$ ranges, and a
per-cell ``predicted vs.\ measured'' validation that confirms
Eq.~\ref{eq:lat-fit} reproduces the median-cell measurement to within
single-digit percent for all measured cells. Predicted and measured
medians agree to within $\sim\!1\%$ on \textsc{MemSearch}-$0_6$B
($849$\,ms predicted vs.\ $890$\,ms measured), giving us empirical
confidence in the composition.

\subsection{Extrapolation across readers for cache-replay backends}
\label{app:lat:extrapolation}

For \textsc{MemSearch} and \textsc{Memobase} we run the answer phase on
\texttt{0\_6b} only. Predictions for heavier readers substitute the
\texttt{0\_6b} prompt-length distribution
$(\widehat{N}_p^{(0\_6b)}, \widehat{N}_c^{(0\_6b)})$ into
Eq.~\ref{eq:lat-fit} with that reader's $(\alpha_r,\beta_r,\gamma_r)$.
This is exact for \textsc{MemSearch} (its retrieved snippets do not
depend on the reader) and approximate for \textsc{Memobase} (its
extractor-specific cache yields slightly different memory snippets per
reader; the prompt format and snippet count are nevertheless stable).
Cells that use this extrapolation are flagged with
\texttt{prompt\_extrapolated\_from\_0\_6b: true} in the output JSON.

\paragraph{Reading Table~\ref{tab:results-L}.}
Each cell decomposes as $T_{\text{search}} + T_{\text{TTFT,LLM}}$:
search overhead from the backend's retrieval pipeline plus the reader's
prefill cost on the resulting prompt. Roman entries are directly
measured; italic entries marked with $\dagger$ have their
$T_{\text{TTFT,LLM}}$ composed from this reader's TTFT fit
(Eq.~\ref{eq:lat-fit-ttft}) applied to \texttt{0\_6b}'s prompt-length
distribution, by the extrapolation rule above. The two bottom rows
report per-backend constants: $T_{\text{search}}$ is the per-question
retrieval median used in every cell of that column, and $N_p$ is the
p50 prompt-token count fed to the reader (taken from the \texttt{0\_6b}
runs, since prompt length is essentially reader-independent for these
five backends). $N_p$ explains why
\textsc{Vanilla} TTFTs sit below \textsc{Oracle}'s at the same reader
even though \textsc{Vanilla} prompts are longer: prefix caching in
\texttt{sglang} reuses the long but identical system prompt across
turns, while \textsc{Oracle}'s shorter but per-question evidence
prompt re-prefills each call. The \texttt{n/a} cell is Mistral-7B,
which we exclude from the latency matrix
(\S\ref{app:lat:caveats}).

\begin{table}[!htbp]
\centering
\small
\caption{Single-stream end-to-end latency $T_{\text{total}}$ on Spark GB10 (ms, p50, \texttt{concurrency=1}). \emph{Roman}: measured cell median (answer-time median plus $T_{\text{search}}$ from Table~\ref{tab:appendix-latency-fit}'s search row). \emph{Italic with $\dagger$}: predicted by composing the per-reader fit (Table~\ref{tab:appendix-latency-fit}) with the \texttt{0\_6b}-measured prompt-length distribution and the backend's $T_{\text{search}}$, where we have not run the cell directly (\S\ref{app:lat:extrapolation}). \texttt{n/a} marks cells with no fit available (Mistral-7B; see \S\ref{app:lat:caveats}).}
\label{tab:appendix-latency-predicted}
\begin{tabular}{lrrrrr}
\toprule
Reader & Vanilla & Oracle & RAG & Memobase & MemSearch \\
\midrule
Qwen3-0.6B & 381 & 279 & 716 & 464 & 890 \\
Llama-3.2-3B & 1215 & 1170 & 1730 & \textit{1800}$^\dagger$ & \textit{2847}$^\dagger$ \\
Mistral-7B-Inst.\ v0.3 & \texttt{n/a} & \texttt{n/a} & \texttt{n/a} & \texttt{n/a} & \texttt{n/a} \\
Qwen3-8B & 2382 & 2486 & 4748 & \textit{3620}$^\dagger$ & \textit{5729}$^\dagger$ \\
Qwen3-32B-AWQ & 5365 & 4380 & 7250 & \textit{5018}$^\dagger$ & \textit{9121}$^\dagger$ \\
\midrule
$T_{\text{search}}$ (ms, p50) & 0 & 0 & 87 & 7 & 48 \\
\bottomrule
\end{tabular}
\end{table}


\subsection{Reproduction recipe}
\label{app:lat:repro}

The full pipeline is, in order:
\begin{enumerate}
\setlength{\itemsep}{0pt}
\item \texttt{scripts/select\_latency\_queries.py}\,$\to$\,\texttt{out/latency\_selection/queries\_s2.json}
\item Phase 1 sweep:
  \texttt{bash scripts/run\_latency\_answer\_sweep.sh} \\(\texttt{READERS=("0\_6b" "llama3b" "7b" "8b" "32b")}, \texttt{BACKENDS=("vanilla" "oracle")})
\item Phase 2.1 (memsearch):
  \texttt{MODEL=0\_6b bash scripts/run\_latency\_memsearch\_ingest.sh}
\item Phase 2.2 (memobase):
  \texttt{MODEL=0\_6b bash scripts/run\_latency\_memobase\_ingest.sh}
\item Fit \& predict: \texttt{python3 scripts/fit\_latency\_model.py}\,$\to$\,\texttt{out/latency\_selection/llm\_fit\_models.json}
\end{enumerate}
Outputs of interest are \texttt{latency\_summary.txt} per cell (p50/p95
ttft, answer time, prompt/completion tokens), \texttt{hw\_agg\_*.json}
per cell (mean/peak GPU power, util, energy/query, peak temp, SoC
rails), and \texttt{llm\_fit\_models.json} (per-reader fit, search
overhead per backend, validation diagnostics, full predicted $5\!\times\!5$
table). Concurrency, mem-fraction, max-running-requests, and the
$8$~egos $\times\,15$~days ingest scope are pinned in the scripts; only
\texttt{MODEL} and \texttt{TRIAL} are intended user-facing knobs.

\subsection{Caveats and limitations}
\label{app:lat:caveats}

\begin{itemize}
\setlength{\itemsep}{0pt}
\item \textbf{Single-stream operating point.} All numbers are
$\text{concurrency}=1$ on both client and server. Throughput-mode
batching would change the constants in Eq.~\ref{eq:lat-fit} significantly
and is not what an edge user experiences.
\item \textbf{Prefix caching.} \texttt{sglang} caches identical prompt
prefixes across requests. \textsc{Vanilla} cells, whose system prompt is
fixed across all $90$ questions, therefore see uncached prefill only on
the per-question evidence tail; this is a deployment-realistic effect
(an edge assistant pinned to a single user reuses prefixes across
turns) and explains why \textsc{Vanilla} TTFTs in Table~\ref{tab:results-L}
sit below \textsc{Oracle} TTFTs at the same reader despite longer raw
prompts.
\item \textbf{8-ego ingest scope.} Phase-2 caches cover the first
$8$~egos. Energy and time scale linearly in egos for both backends
(verified day-by-day in the ingest log), so reporting at $50$~egos is a
straight $\times 50/8$ multiplication where it matters.
\item \textbf{Cache freshness.} The \textsc{MemSearch} cache is
re-built fresh in this work; the \textsc{Memobase} cache reuses an
identically-configured ingest from a prior run on the same node since a
fresh memobase ingest costs $\sim\!72$\,min on \texttt{0\_6b} (extractor
LLM-bound) and the timing is reproducible to within $<\!1\%$.
\item \textbf{Reader-extrapolation for memobase.} As discussed in
\S\ref{app:lat:extrapolation}, \textsc{Memobase} predictions on
non-\texttt{0\_6b} readers assume the prompt-length distribution
transfers from the \texttt{0\_6b} cache.
\item \textbf{Mistral-7B reader.} \texttt{Mistral-7B-Instruct-v0.3}
serves correctly on commodity x86/Hopper deployments but, with the
\texttt{lmsysorg/sglang:spark} build available at submission time on
NVIDIA GB10, exhibits chat-template-induced output collapse: plain
text completion produces coherent prose, while any prompt routed
through \texttt{[INST]\,\dots\,[/INST]} (either via
\texttt{/v1/chat/completions} or as a raw completion) degenerates into
repetitive special-token salad. We verified the weights themselves are
intact (every shard reloads via \texttt{safetensors.safe\_open}; plain
completion is fluent) and that the same failure persists across
\texttt{--chat-template mistral}, the model's bundled chat template,
and across both the \texttt{flashinfer} and \texttt{triton} attention
backends; this matches a known cluster of GB10 / Blackwell / Mistral
incompatibilities under active fix in \texttt{sgl-project/sglang}
(e.g.\ the Qwen3-VL native-mRoPE patch, FA3 / \texttt{sgl\_kernel}
import errors, and SM\_120 RMSNorm-kernel issues). For this submission
we therefore mark the four \texttt{Mistral-7B} cells in
Table~\ref{tab:appendix-latency-predicted} as \texttt{n/a} rather than
extrapolating from a (\textsc{Vanilla}, \textsc{Oracle},
\textsc{RAG}) fit we cannot validate. The methodology and the
single-stream protocol carry over unchanged once the upstream fix
lands.
\end{itemize}
